\newpage
\subsubsection{Herramientas para la aplicación móvil}

La aplicación móvil se desarrollará de forma nativa para el sistema
operativo Android, utilizando un conjunto de herramientas modernas
que garantizan rendimiento, mantenibilidad y compatibilidad con el
enfoque offline-first del sistema.

\textbf{Kotlin} es el lenguaje de programación oficial para el
desarrollo de aplicaciones Android, respaldado por Google desde 2017.
Su sintaxis concisa, su interoperabilidad con Java 21 y su soporte nativo
para corrutinas lo convierten en la opción más adecuada para el
desarrollo de aplicaciones móviles modernas que requieren manejo
eficiente de operaciones asíncronas como la sincronización de datos
y la inferencia del modelo.

\textbf{Jetpack Compose} es el framework moderno de Google para la
construcción de interfaces de usuario en Android. A diferencia del
enfoque tradicional basado en XML, Jetpack Compose permite definir
la interfaz de forma declarativa en Kotlin, reduciendo la cantidad
de código necesario y facilitando la creación de interfaces reactivas
que respondan al estado de la aplicación en tiempo real.

\textbf{TensorFlow Lite} se integrará directamente en la aplicación
para ejecutar el modelo de clasificación EfficientNet-B3 de forma
local en el dispositivo. El modelo en formato \texttt{.tflite} se
almacenará en el dispositivo y se actualizará ocasionalmente desde
Amazon S3 cuando exista una nueva versión disponible.

\textbf{Room} es la biblioteca de persistencia oficial de Android,
construida sobre SQLite. Se utilizará para gestionar la base de datos
local del dispositivo, almacenando los registros de árboles, capturas,
resultados de clasificación y estados de sincronización de forma
estructurada y eficiente mientras el dispositivo opera sin conexión
a internet.

\textbf{MapLibre GL Android} es la biblioteca de código abierto utilizada
para la visualización de mapas vectoriales dentro de la aplicación.
Permite renderizar mapas de forma completamente offline a partir de
un archivo \texttt{.pmtiles} descargado localmente en el dispositivo,
el cual contiene los datos cartográficos de la zona de operación.
Sobre este mapa se renderizan pines de ubicación generados a partir
de las coordenadas GPS extraídas de los metadatos EXIF de cada imagen
capturada, sin requerir conexión a internet en ningún momento durante
su uso.

\textbf{Retrofit y OkHttp} conforman el cliente HTTP que gestionará
la comunicación entre la aplicación móvil y la API REST en la nube
durante la sincronización de datos. Retrofit facilita la definición
de endpoints REST de forma declarativa, mientras que OkHttp actúa
como cliente HTTP subyacente, gestionando las conexiones, los tiempos
de espera y los reintentos automáticos.

La autenticación del experto de campo se gestionará mediante
\textbf{JWT propio}, generado por el backend al momento del inicio
de sesión y almacenado localmente en el dispositivo. Este token se
incluirá en cada petición de sincronización, permitiendo que la
aplicación opere de forma segura sin depender de una conexión
permanente a internet.